\title{Direct Preference Optimization: Your Language Model is Secretly a Reward Model}

\begin{abstract}
Large-scale unsupervised language models (LMs) are capable of acquiring extensive world knowledge and certain reasoning abilities; however, precisely guiding their outputs remains challenging, given the fully unsupervised approach to their training. Traditional methods for making LMs more controllable typically involve obtaining human judgments about the relative merit of different outputs, then fine-tuning the original unsupervised LM to match these preferences—commonly by employing reinforcement learning from human feedback (RLHF). Nevertheless, RLHF introduces significant complexity and potential instability: it requires constructing a reward model that encodes human preferences and subsequently uses reinforcement learning to optimize the LM’s outputs towards this reward, all while avoiding substantial divergence from the base model. In this work, we propose an alternative parameterization of the reward model for RLHF that allows for analytic derivation of the optimal policy, enabling the RLHF problem to be addressed through a straightforward classification objective. This approach, named \textit{Direct Preference Optimization} (DPO), is efficient, robust, and computationally inexpensive, removing the necessity for language model sampling during fine-tuning as well as extensive hyperparameter searches. Empirical results indicate that DPO can match or surpass prior approaches in aligning language models with human preference data. Specifically, DPO-based fine-tuning outperforms PPO-based RLHF in modulating generative sentiment, and attains equal or better performance in tasks such as summarization and single-turn dialogue, all while substantially reducing implementation and training complexity.
\end{abstract}

\section{Introduction}
Large unsupervised language models (LMs) trained on massive datasets have demonstrated unexpected and diverse capabilities~\citep{chowdhery2022palm, brown2020language, touvron2023llama,bubeck2023sparks}. Nonetheless, the training data for these models consists of human-produced content, encompassing a vast range of intentions, competencies, and priorities. Not all such goals or expertise are ideal for the model to replicate. For example, although an AI coding assistant should \textit{recognize} frequent programming errors to better rectify them, we would prefer that it preferentially generate code reflecting the (often uncommon) superior coding skills present within the data. In a similar vein, while it is important for a language model to \textit{comprehend} a widely held misconception (even one believed by 50\% of people), we do not intend for it to affirm that misconception as truth in half of the related interactions. Thus, determining the \emph{intended responses and conduct} of the model, drawing from its broad \textit{reservoir of knowledge and skills}, is a key step toward creating AI systems that are safe, high-performing, and manageable \citep{ouyang2022training}. Although current strategies generally guide LMs to align with human preferences using reinforcement learning (RL), we will demonstrate that the prevailing RL-based objectives can actually be solved directly via a straightforward binary cross-entropy loss, streamlining the preference alignment workflow substantially.

Broadly speaking, contemporary approaches endow language models with preferred behaviors by leveraging curated datasets of human preferences that encode what is perceived as safe and useful behavior. This stage of learning preferences follows a preliminary phase in which the model is exposed to extensive unsupervised pre-training on large text corpora. Supervised fine-tuning on human-generated high-quality responses is the simplest technique for teaching preferences, yet the leading family of approaches utilizes reinforcement learning from human or AI-generated feedback (RLHF/RLAIF; \citep{christiano2017deep,bai2022constitutional}). Within RLHF, a reward model is learned from datasets reflecting human choices, and this reward is then maximized by adjusting the language model's policy through RL, ensuring that its responses achieve high reward scores without substantial deviation from the original model's outputs. Although RLHF yields models with notable prowess in both dialogue and programming, it introduces significant complexity compared to supervised methods, as it necessitates the training of multiple models and on-the-fly sampling during training, resulting in elevated computational demands.

In this work, we present a method to directly train a language model according to human preferences, eliminating the need for explicit reward models or reinforcement learning. We introduce \textit{{Direct Preference Optimization} (DPO)}, a novel approach that implicitly targets the same objective as conventional RLHF strategies—reward maximization subject to a KL-divergence regularization—but offers a more accessible implementation and training procedure. Conceptually, DPO modifies the relative log-likelihoods of preferred versus non-preferred outputs, incorporating an adaptive, instance-specific importance weighting scheme. This mechanism addresses the model collapse observed when using a straightforward log-probability ratio objective. Similar to earlier algorithms, DPO leverages a formal preference framework (for example, the Bradley-Terry model; \cite{bradley1952rankanalysis}) to quantify the congruence between a reward function and actual preference annotations. In contrast to prior approaches, which use this framework to generate a preference loss for reward model training—subsequently followed by policy optimization relative to that reward model—DPO employs a reparameterization that expresses the preference loss as a direct function of the policy itself. With access to human-annotated preferences over generated responses, DPO can train a policy using a straightforward binary cross entropy criterion, ultimately yielding a policy that is optimal with respect to an implicit reward aligned to the preference data.

The principal contribution of our paper is Direct Preference Optimization (DPO), an efficient algorithm for preference-based language model training that avoids reinforcement learning. Through a series of empirical studies, we demonstrate that DPO performs on par with established methods, including PPO-based RLHF, for tasks such as sentiment control, summarization, and conversational response generation, across language models scaling up to 6B parameters.

\section{Related Work}

Self-supervised language models, when scaled up, demonstrate an increasing capacity to handle certain tasks in a zero-shot manner \citep{radford2019language} or through the use of few-shot prompting \citep{gpt3,megatron,chowdhery2022palm}. Nevertheless, their effectiveness on downstream applications and alignment with user objectives are substantially enhanced when fine-tuned on curated datasets comprising instructions and human-authored responses \citep{mishra-etal-2022-cross,sanh2022multitask,chung2022scaling,thoppilan2022lamda}. This process, known as `instruction-tuning,' empowers large language models to better generalize to unseen instructions beyond the initial training distribution, thereby bolstering their practical usability \citep{chung2022scaling}. Although instruction tuning has yielded notable improvements, it is often more feasible to gather \textit{relative} assessments of response quality from non-experts than to amass expert-annotated examples. Consequently, later research has adapted LLMs using datasets informed by human preference rankings, leading to gains in domains such as translation \citep{kreutzer-etal-2018-reliability}, summarization \citep{stiennon2022learning,ziegler2020finetuning}, story generation \citep{ziegler2020finetuning}, and instruction adherence \citep{ouyang2022training,ramamurthy2023is}. These techniques first train a neural reward model to predict preference outcomes consistent with observed preference data, often framed by models like the Bradley-Terry formulation \citep{bradley1952rankanalysis}, and subsequently refine the language model’s outputs through reinforcement learning algorithms such as REINFORCE \citep{williams1992reinforce}, proximal policy optimization (PPO; \cite{schulman2017proximal}), or related adaptations \citep{ramamurthy2023is}. Parallel developments utilize instruction-following LLMs, fine-tuned via human feedback, to synthesize new preference datasets targeting criteria like safety or harmlessness \citep{bai2022constitutional}, using weak supervision supplied only as textual rubrics for annotation. These strategies synthesize progress from two separate research streams: one focused on training language models for diverse objectives through reinforcement learning~\citep{Ranzato2015SequenceLT,paulus2018a,wu2018learning}, and another devoted to developing general-purpose approaches for learning from human preference signals \citep{christiano2017deep,kupcsik2018learning}. Despite the promise inherent in leveraging comparative human feedback, optimizing large language models through reinforcement learning methods is still a significant practical hurdle; this work introduces a theoretically motivated alternative for directly optimizing relative preference data without resorting to RL.

Beyond the domain of language, the problem of learning policies from preferences has been addressed within both bandit and reinforcement learning paradigms, leading to the development of a variety of methodologies. In contextual bandit frameworks, learning may be based on action preferences or rankings in lieu of explicit reward signals; such problems are framed as contextual dueling bandits (CDBs; \cite{yue2012karmed,dudik2015contextual}). In scenarios lacking absolute reward information, the theoretical foundation for CDBs replaces the standard optimal policy with the concept of a \textit{von Neumann winner}, defined as a policy whose expected victory rate over any alternative policy is at least 50\% \citep{dudik2015contextual}. Unlike CDBs, where preferences are obtained in an online fashion, preference-based policy learning from human judgments frequently relies on a predetermined, offline collection of action pairs annotated with preferences \citep{yan2022human}. In a similar vein, \textit{preference-based RL} (PbRL) addresses settings where binary preference feedback, determined by an unobserved scoring function, substitutes for reward feedback \citep{BusaFekete2014,ruiz2023dueling}. Several PbRL algorithms exist, some of which are capable of exploiting off-policy preference datasets; however, these commonly require explicit inference of the underlying scoring or reward function prior to policy optimization \citep{jain2013learning,BusaFekete2014,christiano2017deep,sadigh2017active,kupcsik2018learning}. Contrasting these approaches, we introduce a direct, single-stage policy optimization method that aims to fulfill expressed preferences without separately modeling rewards.

\section{Preliminaries}\label{section:prelims}

Here, we provide an overview of the RLHF framework as articulated in \citeauthor{ziegler2020finetuning} and extended by subsequent works \citep{stiennon2022learning, bai2022training, ouyang2022training}. This process can be delineated into three primary stages: (1) supervised fine-tuning (SFT); (2) collection of preference data and learning a reward function; and (3) policy optimization through reinforcement learning.

\textbf{SFT}: The RLHF protocol typically initiates with the fine-tuning of a pre-trained language model on curated, task-relevant datasets using supervised learning, thereby yielding a model denoted by $\pi^\text{SFT}$.

\textbf{Reward Modelling Phase}: During the second stage, the supervised fine-tuned model is conditioned on prompts $x$ and generates answer pairs $(y_1, y_2) \sim \pi^\text{SFT}(y \mid x)$. Human evaluators are subsequently tasked with indicating their preferences between the two completions, summarized as $y_w\succ y_l \mid x$, with $y_w$ and $y_l$ representing, respectively, the favored and less favored responses among $(y_1, y_2)$. These preference labels are presumed to arise from an unobserved reward function $r^*(y, x)$ to which we have no direct access. Preference data may be modeled in several ways; the Bradley-Terry (BT) model \cite{bradley1952rankanalysis} is a commonly adopted method (though more expressive models such as the Plackett-Luce ranking framework \citep{plackett1975analysis, luce2012individual} are also applicable when more detailed ranked annotations are available). The BT formulation posits that the empirical distribution over human preferences $p^*$ takes the form:

\begin{equation}\label{eq:bradley-terry}
    p^*(y_1\succ y_2 \mid x)=\frac{\exp\left(r^*(x, y_1)\right)}{\exp\left(r^*(x, y_1)\right) + \exp\left(r^*(x, y_2)\right)}.
\end{equation}

Given a fixed dataset of comparison examples $\mathcal{D}=\bigl\{x^{(i)}, y_w^{(i)}, y_l^{(i)}\bigr\}_{i=1}^N$ drawn according to $p^*$, one can introduce a parameterized reward function $r_{\phi}(x, y)$ and infer its parameters through maximum likelihood estimation. By reformulating this as a binary classification task, the negative log-likelihood objective becomes:

\begin{equation}\label{eq:reward_model}
    \mathcal{L}_R(r_{\phi}, \mathcal{D}) = -\mathbb{E}_{(x, y_w, y_l)\sim \mathcal{D}}\bigl[\log \sigma(r_{\phi}(x, y_w)- r_{\phi}(x, y_l))\bigr]
\end{equation}

where $\sigma$ denotes the logistic sigmoid. For language models, $r_{\phi}(x, y)$ typically leverages weights from an SFT model $\pi^\text{SFT}(y \mid x)$ and adds a linear layer on top of the final transformer output to yield a scalar reward \cite{ziegler2020finetuning}. To promote stability and reduce variance in reward estimates, earlier approaches commonly normalize rewards to ensure $\mathbb{E}_{x,y\sim \mathcal{D}}\left[r_\phi(x, y)\right] = 0$ for each $x$.

\textbf{RL Fine-Tuning Phase}: In this stage, the inferred reward function provides the supervision signal for optimizing the language model policy. In line with previous research~\citep{jaques2017sequence, jaques2020human}, policy learning can be expressed as

\begin{equation}\label{eq:RL}
\max_{\pi_{\theta}}  \mathbb{E}_{x\sim \mathcal{D}, y\sim \pi_{\theta}(y \mid x)}\bigl[r_{\phi}(x, y)\bigr] - \beta\mathbb{D}_{\textrm{KL}}\bigl[\pi_{\theta}(y\mid x)\mid \mid \pi_\text{ref}(y\mid x)\bigr],
\end{equation}

with $\beta$ governing the regularization against deviation from a reference policy $\pi_\text{ref}$, typically set to the initial SFT policy $\pi^\text{SFT}$. The language model $\pi_\theta$ is initialized from $\pi^\text{SFT}$ as well. This KL constraint plays a crucial role, as it limits divergence from domains where the reward model remains valid and discourages the policy from collapsing into narrow regions of high reward, thereby preserving response diversity. Since text generation is inherently discrete and non-differentiable, this optimization is handled via reinforcement learning methods. The prevailing strategy \citep{ziegler2020finetuning, stiennon2022learning, bai2022training, ouyang2022training} formulates the reward for RL as ${r(x, y) = r_{\phi}(x, y) -\beta (\log \pi_{\theta}(y\mid x) - \log \pi_\text{ref}(y\mid x))}$, and optimizes it with algorithms such as PPO \cite{schulman2017proximal}.

\section{Direct Preference Optimization}\label{sec:DPO}

Recognizing the difficulties inherent in applying reinforcement learning techniques to large-scale fine-tuning of language models, we set out to formulate a streamlined method for preference-based policy learning. Distinct from traditional RLHF workflows that first learn a reward model before optimizing it via RL, our framework makes use of a special reward parameterization, allowing direct recovery of the optimal policy analytically—bypassing the need for iterative RL optimization. As elaborated in the following sections, the central idea is to exploit a closed-form correspondence between the reward function and its induced optimal policy, enabling us to recast losses defined over rewards as losses over policies. By changing variables in this manner, we avoid explicit reward model training but still fit models consistent with standard preference structures, such as the Bradley-Terry model. In this paradigm, the policy architecture simultaneously encodes both the generative model and an implicit reward signal.

\textbf{Deriving the DPO objective.} Our analysis begins with the reinforcement learning objective found in earlier studies, Eq.~\ref{eq:RL}, which utilizes a general reward function $r$. In line with existing literature~\citep{peters2007reinforcement, peng2019advantage, korbak2022reinforcement, go2023aligning}, one can readily demonstrate that the optimal policy for the KL-regularized reward maximization problem in Eq.~\ref{eq:RL} can be written as:

\begin{equation}\label{eq:op_policy}
    \pi_r(y\mid x) = \frac{1}{Z(x)}\pi_\text{ref}(y\mid x)\exp\left(\frac{1}{\beta}r(x, y)\right),
\end{equation}

with $Z(x) =\sum_{y}\pi_\text{ref}(y\mid x)\exp\left(\frac{1}{\beta}r(x, y)\right)$ denoting the normalization constant, or partition function. Appendix \ref{app:derivation1} presents the detailed derivation. In practical scenarios, even when employing the maximum likelihood estimate $r_\phi$ of the actual reward $r^*$, directly computing the partition function $Z(x)$ remains computationally intensive \citep{korbak2022reinforcement, go2023aligning}, limiting the practicality of this formulation. Nonetheless, Eq.~\ref{eq:op_policy} can be manipulated to solve for the reward function as a function of its associated optimal policy $\pi_r$, the reference policy $\pi_\text{ref}$, and the partition function $Z(\cdot)$. This is achieved by taking the logarithm of Eq.~\ref{eq:op_policy}, and after some algebraic manipulation, we obtain:

\begin{equation}\label{eq:main_eq}
    r(x,y) =\beta \log \frac{\pi_r(y\mid x)}{\pi_\text{ref}(y\mid x)} + \beta \log Z(x).
\end{equation}

Applying this reformulation to the true reward $r^*$ and the corresponding optimal policy $\pi^*$, we note that the Bradley-Terry model relies solely on reward differences between candidate outputs, specifically, ${p^*(y_1 \succ y_2 \mid x) = \sigma(r^*(x, y_1) - r^*(x, y_2))}$. By inserting the reparameterized form for $r^*(x,y)$ from Eq.~\ref{eq:main_eq} into the Bradley-Terry preference model (Eq.~\ref{eq:bradley-terry}), the normalization term cancels out. As a result, the probability of a human preferring one completion over another can be expressed purely in terms of the optimal policy $\pi^*$ and the reference policy $\pi_\text{ref}$. Thus, the optimal RLHF policy $\pi^*$ under the Bradley-Terry framework satisfies the following preference relationship:

\begin{equation}\label{eq:objective}
    p^*(y_1\succ y_2 \mid x)=\frac{1}{1 + \exp\left(\beta \log \frac{\pi^*(y_2\mid x)}{\pi_\text{ref}(y_2\mid x)} - \beta \log \frac{\pi^*(y_1\mid x)}{\pi_\text{ref}(y_1\mid x)}\right)}
\end{equation}

See Appendix~\ref{app:derivation2} for the step-by-step derivation. Although Eq.~\ref{eq:objective} utilizes the Bradley-Terry model, analogous formulations can be constructed for more general Plackett-Luce models~\citep{plackett1975analysis, luce2012individual}; see Appendix~\ref{app:plackett_luce_models} for details.

Having recast the likelihood of observed human preference data in terms of the optimal policy, as opposed to a reward model, we are now equipped to formulate a maximum likelihood estimation objective for a parameterized policy $\pi_\theta$. In analogy with reward model-based objectives (such as Eq.~\ref{eq:reward_model}), this leads to the following policy objective:

\begin{equation}\label{eq:optimum_model}
    \mathcal{L}_\text{DPO}(\pi_{\theta}; \pi_\text{ref}) = -\mathbb{E}_{(x, y_w, y_l)\sim \mathcal{D}}\left[\log \sigma \left(\beta \log \frac{\pi_{\theta}(

In this formulation, an implicit reward is learned through an alternative parameterization such that the resulting optimal policy aligns with $\pi_\theta$. Furthermore, since our approach corresponds to fitting a reparameterized Bradley-Terry model, it inherits several theoretical guarantees, including consistency provided that the preference data satisfies certain conditions \cite{bong2022generalized}. Additional theoretical considerations and their relation to prior research are explored in Section~\ref{sec:theory}.

\textbf{What is the mechanism behind the DPO update?} To gain a more precise intuition about how DPO operates, one can examine the gradient of its loss function $\mathcal{L}_\text{DPO}$. The gradient with respect to the parameter vector $\theta$ takes the following form:

\begin{multline*}\label{eq:gradient}
    \nabla_\theta \mathcal{L}_\text{DPO}(\pi_\theta;\pi_\text{ref}) = \\ -\beta\mathbb{E}_{(x, y_w, y_l) \sim \mathcal{D}} \bigg[\underbrace{\sigma(\hat{r}_\theta(x, y_l) - \hat{r}_\theta (x, y_w))}_\text{higher weight when reward estimate is wrong}\bigg[\underbrace{\nabla_\theta\log \pi(y_w \mid x)}_\text{increase likelihood of $y_w$} - \underbrace{\nabla_\theta\log\pi(y_l \mid x)}_\text{decrease likelihood of $y_l$}\bigg]\bigg],
\end{multline*}

where $\hat{r}_\theta(x, y) = \beta \log \frac{\pi_\theta(y \mid x)}{\pi_\text{ref}(y \mid x)}$ denotes the reward function implicitly determined by the language model $\pi_\theta$ relative to the reference policy $\pi_\text{ref}$ (see Section~\ref{sec:theory} for further elaboration). Conceptually, the gradient of $\mathcal{L}_\text{DPO}$ serves to reinforce the model’s probability of generating favored completions $y_w$, while diminishing the likelihood of less preferred completions $y_l$. Crucially, each example is reweighted according to the degree by which the implicit reward $\hat{r}_\theta$ incorrectly ranks $y_l$ above $y_w$, modulated by $\beta$, thus incorporating both the severity of ranking errors and the influence of the KL-divergence constraint. Our empirical results indicate that this specific reweighting is essential, as a simpler variant of the method omitting the weighting term can lead to model collapse (see Appendix Table~\ref{tab:unlikelihood_generations}).

\textbf{DPO overview.} The standard DPO workflow consists of: 1) Drawing response samples $y_1, y_2 \sim \pi_\text{ref}(\cdot \mid x)$ for each prompt $x$ and annotating them with human preferences to create the offline preference dataset $\mathcal{D} = \{x^{(i)}, y_w^{(i)}, y_l^{(i)}\}_{i=1}^N$; and 2) training the language model $\pi_\theta$ by minimizing $\mathcal{L}_\text{DPO}$, conditioned on the chosen $\pi_\text{ref}$, dataset $\mathcal{D}$, and the target $\beta$. In practical scenarios, preference datasets that are already publicly available are preferable to constructing new samples and manually collecting human preferences. Because these datasets are typically generated using $\pi^\text{SFT}$, we set $\pi_\text{ref} = \pi^\text{SFT}$ whenever this baseline is accessible. Conversely, if $\pi^\text{SFT}$ is unavailable, we obtain $\pi_\text{ref}$ by maximizing the likelihood over the set of preferred completions ${(x, y_w)}$, i.e., ${\pi_\text{ref} = \argmax_{\pi}\mathbb{E}_{x, y_w \sim \mathcal{D}}\left[\log \pi(y_w \mid x)\right]}$. This step reduces the discrepancy between the actual, yet inaccessible, reference distribution and the surrogate $\pi_\text{ref}$ employed in DPO. More implementation specifics and details regarding hyperparameters are included in Appendix~\ref{app:implementation}.

\section{Theoretical Analysis of DPO} This section further elucidates the DPO approach, develops its theoretical justification, and connects the benefits of DPO to challenges found in actor-critic algorithms commonly applied in RLHF, such as PPO~\cite{schulman2017proximal}.

\label{sec:theory}

\subsection{Your Language Model Is Secretly a Reward Model} DPO achieves policy optimization without explicitly modeling a reward function or conducting reinforcement learning, instead using a single maximum likelihood-based objective. Specifically, the loss function in Eq. \ref{eq:main_eq} corresponds to the Bradley-Terry framework with the reward structure $r^*(x, y) = \beta \log\frac{\pi^*_\theta(y \mid x)}{\pi_\text{ref}(y \mid x)}$; optimization of our parameterized model $\pi_{\theta}$ thus mirrors the process of reward model fitting described in Eq. \ref{eq:reward_model} through an appropriate transformation of variables. In this section, we formalize the theory underpinning this reparameterization, demonstrate that it preserves the flexibility of learned reward models, and show that it can exactly recover the optimal policy. We begin by introducing an equivalence relation for reward functions.

\begin{definition}
Two reward functions $r(x, y)$ and $r'(x, y)$ are considered equivalent if and only if there exists a function $f$ such that ${r(x, y)-r'(x, y) = f(x)}$.
\end{definition}

It is straightforward to verify that this defines an equivalence relation, resulting in a partition of the space of reward functions into distinct classes. We next present two relevant lemmas:

\begin{lemma}\label{lemma:same_prefrence} Within the Plackett-Luce framework, including the specific case of Bradley-Terry, any two reward functions belonging to the same equivalence class yield identical distributions over preferences.
\end{lemma}

\begin{lemma}\label{lemma:same_policy}
    For the constrained RL formulation, any pair of reward functions from the same equivalence class will lead to the same optimal policy.
\end{lemma}

The proofs follow directly and can be found in Appendix \ref{app:lemma1}. The first lemma reiterates a widely recognized issue of under-specification in the Plackett-Luce family of models \cite{plackett1975analysis}. Because of this ambiguity, it is typically necessary to enforce extra identifiability constraints to secure guarantees for maximum likelihood estimation in Eq. \ref{eq:reward_model} \cite{bong2022generalized}. The second lemma clarifies that, since all reward functions within a given class produce the same optimal policy, it is sufficient for our purposes to recover any single representative from the optimal equivalence class. In Appendix~\ref{app:thm1}, we establish the following theorem:

\begin{theorem}\label{thm:main}
    Assuming mild regularity conditions, every equivalence class of reward functions compatible with Plackett-Luce (including Bradley-Terry as a special case) can be expressed using the reparameterization ${r(x, y) = \beta \log \frac{\pi(y\mid x)}{\pi_\text{ref}(y\mid x)}}$ for some model $\pi(y\mid x)$ and some specified reference model $\pi_\text{ref}(y \mid x)$.
\end{theorem}

\begin{sproof}
    Let $r(x, y)$ denote an arbitrary reward function that induces an optimal model $\pi_r(y \mid x)$, as defined by Eq. \ref{eq:op_policy}. We aim to show that one can always construct a member of the equivalence class of $r$ using the reparameterization above. Define the projection $f$ by  
\begin{equation}
    f(r; \pi_\text{ref}, \beta)(x, y) = r(x, y) - \beta\log\sum_{y}\pi_\text{ref}(y\mid x)\exp\left(\frac{1}{\beta}r(x, y)\right)
\end{equation}
This projection normalizes the reward function by subtracting the log-partition function for $\pi_r$. Since the normalization term depends solely on $x$, the function $f(r; \pi_\text{ref}, \beta)(x, y)$ remains within the equivalence class of $r(x, y)$. Substituting the right-hand side of Eq.~\ref{eq:main_eq} (which universally characterizes any reward function), we obtain $f(r; \pi_\text{ref}, \beta)(x, y) = \beta \log \frac{\pi_r(y\mid x)}{\pi_\text{ref}(y\mid x)}$. Thus, $f$ yields a form-equivalent representative of $r$ matching the required structure, ensuring our reward model is maximally expressive under this reparameterization.
\end{sproof}

Theorem~\ref{thm:main} can be reinterpreted as identifying the specific reward function, within each equivalence class, that the DPO reparameterization chooses. In particular, it selects the reward function that fulfills:

\begin{equation}\label{eq:lag_p}
     \sum_{y}\underbrace{\pi_\text{ref}(y\mid x)\exp\left(\frac{1}{\beta}r(x, y)\right)}_{=\pi(y\mid x)\text{, using Thm.~\ref{thm:main} reparam.}} = 1,
\end{equation}

which ensures that $\pi(y\mid x)$ defines a legitimate probability distribution (all probabilities are non-negative and sum to 1). Looking at Eq.~\ref{eq:op_policy}, we observe that Eq.~\ref{eq:lag_p} represents the partition function of the optimal policy determined by the reward function $r(x, y)$. The principal contribution of the DPO algorithm is to enforce constraints on the underdetermined Plackett-Luce family of preference models (with a particular emphasis on the Bradley-Terry model). This constraining maintains the expressivity of the reward models but, importantly, renders the optimal policy in Eq. \ref{eq:op_policy} analytically computable for every prompt $x$.

\subsection{Instability of Actor-Critic Algorithms}
The framework presented can further be utilized to shed light on the instabilities commonly observed in standard actor-critic algorithms employed in RLHF, such as PPO. In line with the RLHF methodology, we concentrate on the RL fine-tuning phase described in Section \ref{section:prelims}. Notably, one can draw parallels with the control as inference approach \cite{levine2018reinforcement} as applied to the constrained RL problem detailed in \ref{eq:RL}. Here, a parameterized policy $\pi_{\theta}(y\mid x)$ is considered, with the aim of minimizing $\mathbb{D}_{\text{KL}}[\pi_{\theta}(y|x) \mid \mid \pi^*(y\mid x)]$, where $\pi^*$ denotes the optimal policy defined in Eq. \ref{eq:optimum_model} and is based on the reward function $r_{\phi}(y, x)$. After rearranging terms, this results in the following optimization criterion:

\begin{equation}\label{eq:AC}
    \max_{\pi_{\theta}}\mathbb{E}_{\pi_{\theta}(y\mid x)}\bigg[\underbrace{r_{\phi}(x, y) -\beta\log\sum_{y}\pi_\text{ref}(y\mid x)\exp\left(\frac{1}{\beta}r_{\phi}(x, y)\right)}_{f(r_{\phi}, \pi_\text{ref}, \beta)} - \underbrace{\beta\log\frac{\pi_{\theta}(y\mid x)}{\pi_\text{ref}(y\mid x)}}_{\text{KL}}\bigg]
\end{equation}

The objective considered here is identical to that optimized in previous studies 
\citep{ziegler2020finetuning, stiennon2022learning, bai2022training, ouyang2022training} for the reward class $r_{\phi}$, where the reward is DPO-equivalent. In this context, the normalization component of $f(r_{\phi}, \pi_\text{ref}, \beta)$ can be seen as representing the soft value function for the reference policy $\pi_\text{ref}$. Although this normalization factor does not influence the final solution, omitting it may result in a high-variance policy gradient, which can destabilize the learning process. One possible strategy is to model the normalization term with a learned value function, yet optimizing such a function poses its own challenges. Previous approaches have sometimes relied on reward normalization via a human completion baseline, effectively approximating the normalizer with a single Monte-Carlo sample. By contrast, the DPO reparameterization provides a reward function formulation that eliminates the need for any explicit baselines.

\section{Experiments}
This section presents an empirical assessment of DPO in training policies directly from preference data. We begin with a controlled text-generation task, examining DPO's efficiency in balancing reward maximization with minimizing KL-divergence relative to the reference policy, and compare its performance to widely-used preference learning approaches like PPO. We then move to evaluating DPO on larger-scale models and more complex RLHF challenges, such as summarization and dialogue. Our results indicate that DPO, with minimal hyperparameter adjustment, typically matches or surpasses robust baselines like PPO-based RLHF, and can outperform returning the best out of $N$ sampled trajectories judged by a learned reward model. Prior to discussing these outcomes, we provide an overview of our experimental design; supplementary methodological details can be found in Appendix~\ref{app:exp_details}.

\textbf{Tasks.} Our empirical study investigates three distinct open-ended text generation tasks. Across all tasks, the methods optimize a policy using a dataset of human preferences $\mathcal{D}=\bigl\{x^{(i)}, y_w^{(i)}, y_l^{(i)}\bigr\}_{i=1}^N$. For the \textbf{controlled sentiment generation} task, the input $x$ consists of the opening segment of a movie review selected from the IMDb dataset \cite{maas-EtAl:2011:ACL-HLT2011}, and the objective is for the policy to generate an output $y$ exhibiting positive sentiment. To facilitate controlled assessments in this scenario, we automatically construct preference pairs based on generations evaluated by a pre-trained sentiment classifier, such that $p(\text{positive}\mid x,y_w)>p(\text{positive}\mid x,y_l)$. The SFT baseline is produced by fine-tuning GPT-2-large on the training partition of the IMDb corpus, continuing the process until model convergence (see App~\ref{app:sentiment_details} for comprehensive information). For the \textbf{summarization} task, $x$ corresponds to a Reddit forum post, and the policy is tasked with composing a summary $y$ encapsulating the main points. In alignment with existing literature, we employ the Reddit TL;DR dataset \citep{volske-etal-2017-tl} as well as the human preference annotations collected by \citeauthor{stiennon2022learning}. The SFT model used here is fine-tuned on summaries authored by humans for forum posts,\footnote{\url{https://huggingface.co/CarperAI/openai_summarize_tldr_sft}} utilizing the TRLX \citep{leandro_von_werra_2023_7790115} toolkit for RLHF training. Human preference data originate from samples produced by a comparable SFT model, as documented by \citeauthor{stiennon2022learning}. The third task, \textbf{single-turn dialogue}, involves $x$ as a prompt posed by a user—ranging from scientific queries to personal advice. Here, the policy is required to provide a response $y$ that is both informative and supportive. This evaluation employs the Anthropic Helpful and Harmless dialogue dataset \citep{bai2022training}, comprising 170,000 exchanges between humans and an AI assistant. Each dialogue session concludes with two model-generated candidate replies, alongside a human label indicating the preferred choice. Due to the lack of a dedicated pre-trained SFT model for this setting, we instead adapt a publicly available language model by fine-tuning it on those completions rated as preferred, thereby constructing the SFT model.

\textbf{Evaluation.} We employ two distinct evaluation methodologies in our experiments. For controlled sentiment generation, we assess the capacity of each algorithm to optimize the constrained reward maximization objective by computing its Pareto frontier between the achieved reward and the KL-divergence from the reference policy. This analysis is feasible as the ground-truth reward function—a sentiment classifier—is available. Conversely, in practical scenarios where the ground truth reward is inaccessible, we measure the performance of algorithms through their \textit{win rate} when compared to a baseline policy. Specifically, in the summarization and single-turn dialogue tasks, we rely on GPT-4 as a stand-in for human judgment to evaluate summary quality and the helpfulness of dialogue responses, respectively. In the summarization setting, the test set’s reference summaries serve as the baseline, while for dialogue, the baseline is the annotator-preferred response from the test data. Despite prior findings suggesting that language models can serve as superior automated evaluators relative to conventional metrics \citep{Chen2023ExploringTU}, we supplement our automatic evaluation with a human study, presented in Sec.~\ref{sec:human-judgments}, to validate our reliance on GPT-4. The results indicate that GPT-4’s assessments are highly aligned with human judgment, with human-GPT-4 agreement typically equaling or surpassing agreement between human annotators.

\textbf{Methods.} Alongside DPO, we benchmark various established methods for aligning language model behavior with human preferences. In the summarization experiments, we apply zero-shot prompting to \textbf{GPT-J} \citep{gpt-j}, and in dialogue, we employ 2-shot prompting with \textbf{Pythia-2.8B} \citep{biderman2023pythia}. Additionally, we test the \textbf{SFT} approach and \textbf{Preferred-FT}, a model obtained by supervised fine-tuning on the chosen output $y_w$—generated either by SFT (for sentiment control and summarization) or a standard language model (for single-turn dialogue). We further investigate the pseudo-supervised \textbf{Unlikelihood} method~\citep{welleck2019neural}, which enhances the likelihood of $y_w$ while explicitly penalizing the model for assigning probability to $y_l$; an adjustable coefficient $\alpha\in[0,1]$ modulates the strength of the unlikelihood component. Our study also includes \textbf{PPO} \citep{schulman2017proximal}, leveraging a reward function trained on preference data, and \textbf{PPO-GT}, an oracle variant trained with access to the true reward function within the controlled sentiment task. For the sentiment experiments, we utilize both an off-the-shelf PPO-GT implementation \cite{leandro_von_werra_2023_7790115} and a custom version that applies reward normalization and improved hyperparameter tuning, modifications that are also adopted for standard PPO when using learned rewards. Lastly, we include the \textbf{Best of $N$} baseline: $N$ candidate responses are generated from the SFT model (or Preferred-FT in the dialogue context), and the highest-rewarding completion—as determined by a reward model trained on preference data—is selected. While this approach is capable of strong performance and effectively separates reward model accuracy from PPO optimization, it is resource-intensive, requiring $N$ samples per test input and thus becomes impractical for even moderately sized $N$.

\subsection{How well can DPO optimize the RLHF objective?}

RLHF algorithms typically aim to maximize reward while penalizing divergence from a reference policy through a KL constraint, balancing the incentive to exploit high-reward actions with the necessity to remain close to prior behavior. Consequently, when evaluating and comparing these methods, it is necessary to jointly consider both the achieved reward and the magnitude of KL divergence; increasing reward at the cost of disproportionately higher KL is not inherently beneficial. In Figure~\ref{fig:frontier-tldr-main}, we display the trade-off frontier between reward and KL for different methods within the sentiment task. To probe this, we conduct multiple independent training experiments for each method, adjusting a parameter governing policy conservativeness for each (specifically: target KL $\in\{3,6,9,12\}$ for PPO, $\beta \in \{0.05,0.1,1,5\}$, $\alpha\in\{0.05,0.1,0.5,1\}$ for unlikelihood, and various random seeds for preferred-FT), totaling 22 training runs. At intervals of 100 training updates up to convergence, we assess each policy on a shared set of held-out prompts, measuring both the mean true reward and the average sequence-level KL\footnote{That is, the sum of the per-timestep KL-divergences.} relative to the reference policy, i.e., $\text{KL}\left(\pi\mid \mid \pi_\text{ref}\right)$. Our findings indicate that DPO attains a substantially more optimal frontier, securing the highest rewards at consistently low KL levels. This outcome is significant for several reasons. First, despite optimizing an identical objective to PPO, DPO demonstrates clear superiority in reward/KL efficiency—its trade-off strictly dominates that of PPO. Second, even when PPO is supplied with oracle access to true rewards (PPO-GT), DPO still establishes a superior reward-KL frontier.

\subsection{Can DPO scale to real preference datasets?}
\label{sec:dpo-real-datasets}
We proceed by assessing the fine-tuning capabilities of DPO on tasks involving summarization and single-turn dialogue. In the context of summarization, widely used automated metrics like ROUGE often show weak correlation with human judgment~\citep{stiennon2022learning}. Previous studies have demonstrated that leveraging PPO-based fine-tuning on human preference data can yield superior summaries. To benchmark the different approaches, we generate model outputs on the TL;DR summarization dataset's test split and calculate the average win rate of these completions against references within the test set. Completions are sampled across a range of temperatures from 0.0 to 1.0, and the corresponding win rates are illustrated in Figure~\ref{fig:frontier-tldr-main} (right panel). For comparability, DPO, PPO, and Preferred-FT all start from the identical GPT-J SFT model\footnote{\url{https://huggingface.co/CarperAI/openai_summarize_tldr_sft}}. Notably, DPO achieves an approximate 61\% win rate at a sampling temperature of 0.0, outperforming PPO, which attains around 57\% at its own optimal temperature setting. DPO also surpasses the highest win rate achieved by the best-of-$N$ baseline. It is important to mention that minimal effort was directed towards optimizing DPO's $\beta$ parameter, thus the reported outcomes might not fully reflect DPO's maximal effectiveness. Additionally, DPO exhibits much greater resilience to changes in sampling temperature compared to PPO, whose performance can deteriorate at higher temperatures to the level of the unmodified GPT-J baseline. In contrast, Preferred-FT does not demonstrate notable improvements relative to the SFT baseline. A direct human evaluation between DPO and PPO, detailed in Section~\ref{sec:human-judgments}, indicates that human raters favored DPO samples at temperature 0.25 in 58\% of cases over PPO samples at the same temperature.

For single-turn dialogue evaluation, we assess the various approaches on a subset of the Anthropic HH dataset test split \citep{bai2022training}, specifically focusing on cases involving a single exchange between a human and an assistant. For GPT-4 assessments, the model's preferred completions from the test set serve as the ground truth for determining the win rates of each method. In the absence of a widely accepted SFT baseline, we initiate with the pre-trained Pythia-2.8B, apply Preferred-FT to fine-tune a reference model on the selected completions to ensure the outputs match the distribution, and subsequently train with DPO. Additionally, we benchmark against the strongest result among 128 Preferred-FT completions (having observed a performance plateau at $N=128$ for this scenario; see Appendix Figure~\ref{fig:best-of-n}) and compare to a 2-shot prompt approach with the Pythia-2.8B base model. Our analysis reveals that DPO achieves performance on par with or superior to the alternatives when optimal temperature settings are used for each method. We also include an RLHF baseline, employing PPO and trained on the Anthropic HH dataset \footnote{\url{https://huggingface.co/reciprocate/ppo_hh_pythia-6B}}, implemented using a widely-recognized codebase \footnote{\url{https://github.com/CarperAI/trlx/tree/main/examples/hh}}, but we are unable to identify prompt or sampling temperature configurations that outperform the base Pythia-2.8B. Given our TL;DR findings and the shared reward optimization goal between the approaches, we treat the Best of 128 benchmark as an approximate indicator for PPO-level outcomes. Overall, DPO emerges as the only method that, while maintaining computational efficiency, surpasses the preferred completions in the Anthropic HH dataset and matches or exceeds the results of the resource-intensive Best of 128 baseline. Lastly, Figure~\ref{fig:dialogue-main} demonstrates that DPO attains peak performance after relatively few training steps.

\subsection{Generalization to a new input distribution}

To better understand how PPO and DPO handle distributional changes, we assess the policies trained on the Reddit TL;DR summarization task using data from a distinct source: news articles in the test split of the CNN/DailyMail dataset \citep{nallapati-etal-2016-abstractive}. The policies are evaluated with the same optimal sampling temperatures identified on TL;DR (0 and 0.25). The comparative results can be found in Table~\ref{tab:ood}. To determine performance, we calculate the win rate of GPT-4 relative to the ground-truth summaries for each dataset, adapting the original GPT-4 (C) evaluation prompt from Reddit TL;DR by substituting “forum post” with “news article”. On this new domain, DPO still exhibits a marked performance advantage over PPO. These findings offer preliminary support for the notion that DPO-trained policies generalize to unfamiliar distributions as effectively as those from PPO, even though DPO does not make use of the supplementary unlabeled Reddit TL;DR prompts incorporated into PPO’s training.

\subsection{Validating GPT-4 judgments with human judgments}
\label{sec:human-judgments}
We undertake a human evaluation to corroborate the trustworthiness of GPT-4's assessments, focusing on outputs from the TL;DR summarization task and employing two distinct GPT-4 prompts. The \textbf{GPT-4 (S)} (simple) prompt only asks which summary better captures the key information of the post, whereas the \textbf{GPT-4 (C)} (concise) prompt adds a requirement for conciseness; the latter is examined as we observe GPT-4 favoring lengthier, more repetitive summaries when using the former. Appendix~\ref{app:prompts} provides the full prompt texts. Three specific method pairings are compared—one with the strongest-performing system (DPO at temperature 0.25), one with the weakest (PPO at temperature 1.0), and a middle case (SFT at temperature 0.25)—to cover a range of output qualities. Each method is benchmarked against PPO with greedy sampling (its optimal temperature setting). Results show that GPT-4’s preferences, according to both prompts, align with human judgments at rates comparable to the agreement among human annotators themselves; due to the restricted pool of raters, multiple human annotations are collected only for the DPO and PPO-1 comparisons. In general, the \textbf{GPT-4 (C)} prompt yields outcomes that best reflect human preferences, which motivates its use for reporting primary results in Section~\ref{sec:dpo-real-datasets}. More information regarding the human evaluation procedure, including the rater-facing web interface and roster of participating annotators, is available in Appendix~\ref{app:human-study}.

\section{Discussion}
Preference-based learning stands out as a robust and scalable method for developing language models that are both capable and aligned with user expectations. In this work, we presented DPO, a streamlined approach that allows language models to be trained directly from preference data, circumventing the need for reinforcement learning. Unlike methods that reframe preference learning as a reinforcement learning (RL) challenge in order to leverage standard RL tools, DPO establishes a principled connection between policy learning in language models and associated reward functions. This relationship facilitates optimization toward human preferences using a straightforward cross-entropy objective, entirely avoiding reinforcement learning without any sacrifice of generality. DPO achieves comparable or superior results to existing RLHF strategies, including those utilizing PPO, and does so with negligible hyperparameter tuning—substantially lowering the barriers to utilizing human preferences for language model training.

\textbf{Limitations \& Future Work.} Several important avenues for further investigation are highlighted by our findings. One question is the extent to which policies learned via DPO can generalize to distributions beyond the training data, relative to approaches that rely on explicit reward functions. Preliminary experiments indicate that DPO-trained policies show generalization performance on par with PPO-trained counterparts, yet this requires a more thorough assessment. For instance, there is potential to examine whether using the DPO model for self-labeling can help leverage unlabeled prompts effectively. Additionally, further analysis is warranted regarding how reward over-optimization might appear in the direct preference optimization context—specifically, whether the modest reduction in performance observed in Figure~\ref{fig:dialogue-main}-right is indicative of this phenomenon. While our evaluations included models with up to 6B parameters, extending DPO to much larger, state-of-the-art models presents a compelling prospect for upcoming work. With respect to evaluation, we observe that win rates as assessed by GPT-4 are sensitive to prompt phrasing; identifying optimal methods for obtaining reliable automated evaluations represents an interesting future direction. Finally, DPO’s utility is not confined to language model alignment via human preferences—its principles may also be applied to the training of generative models in a variety of other domains.